\section{Experiments}
\label{sec:experiments}

We evaluate TarMAC on a variety of tasks and environments.
All our models were trained with a batched synchronous version of the multi-agent
Actor-Critic described above, using RMSProp with a learning rate of
$7\times 10^{-4}$ and $\alpha = 0.99$, batch size $16$, discount factor
$\gamma = 0.99$ and entropy regularization coefficient $0.01$ for agent policies.
All our agent policies are instantiated from the same set of shared parameters; \ie $\theta_1 = ... = \theta_N$.
Each agent's GRU hidden state is $128$-d, message signature/query is $16$-d, and message value is $32$-d (unless specified otherwise).
All results are averaged over $5$ independent seeds (unless noted otherwise),
and error bars show standard error of means.

\vspace{-5pt}
\subsection{SHAPES}
\label{sec:shapes}

The SHAPES dataset was introduced by~\citet{andreas_cvpr16}\footnote{\href{https://github.com/jacobandreas/nmn2/tree/shapes}{\texttt{github.com/jacobandreas/nmn2/tree/shapes}}}, and originally created
for testing compositional visual reasoning for the task of visual question answering.
It consists of synthetic images of $2$D colored shapes
arranged in a grid ($3 \times 3$ cells in the original dataset) along with corresponding
question-answer pairs.
There are $3$ shapes (circle, square, triangle), $3$ colors (red, green, blue),
and $2$ sizes (small, big) in total (see~\Figref{fig:shapes_3x3}).

We convert each image from SHAPES
into an active environment where agents can now be spawned at different regions of
the image, observe a $5 \times 5$ local patch around them and their coordinates, and take actions
to move around -- $\{$\texttt{up}, \texttt{down}, \texttt{left}, \texttt{right}, \texttt{stay}$\}$. Each agent is tasked with
navigating to a specified goal state in the environment within a max no.~of steps -- $\{$`red', `blue square',
`small green circle', \etc$\}$ -- and the reward for each agent at every timestep
is based on team performance \ie $r_t = \frac{\textnormal{\# agents on goal}}{\textnormal{\# agents}}$.
Having a symmetric, team-based reward incentivizes agents to cooperate
in finding each agent's goal.

\begin{table*}[h!]
\centering
\begin{tabular}{lrrr}
   & \multicolumn{1}{c}{{\scriptsize $30 \times 30$, $4$ agents, find$[$\texttt{red}$]$}} & \multicolumn{1}{c}{{\scriptsize $50 \times 50$, $4$ agents, find$[$\texttt{red}$]$}} & \multicolumn{1}{c}{{\scriptsize $50 \times 50$, $4$ agents, find$[$\texttt{red,red,green,blue}$]$}}    \\ \midrule
    {\small No communication}
        & {\small $95.3$}{\scriptsize $\pm 2.8\%$} & {\small $83.6$}{\scriptsize $\pm 3.3\%$} & {\small $69.1$}{\scriptsize $\pm 4.6\%$}  \\[0.05in]
    {\small No attention}
        & {\small $\mathbf{99.7}$}{\scriptsize $\pm 0.8\%$} & {\small $\mathbf{89.5}$}{\scriptsize $\pm 1.4\%$} & {\small $82.4$}{\scriptsize $\pm 2.1\%$} \\[0.05in]
    {\small TarMAC}
        & {\small $\mathbf{99.8}$}{\scriptsize $\pm 0.9\%$} & {\small $\mathbf{89.5}$}{\scriptsize $\pm 1.7\%$} & {\small $\mathbf{85.8}$}{\scriptsize $\pm 2.5\%$}  \\ \bottomrule
\end{tabular} \\[0.05in]
\caption{Success rates on $3$ different settings of cooperative navigation
    in the SHAPES environment.}
\label{table:shapesnumbers}
\vspace{-15pt}
\end{table*}

\textbf{How does targeting work?}
Recall that each agent predicts a signature and value vector as part of the message
it sends, and a query vector to attend to incoming messages.
The communication is targeted because the attention probabilities are a function
of both the sender's signature and receiver's query vectors. So it is not just
the receiver deciding how much of each message to listen to.
The sender also sends out signatures that affects how much of each
message is sent to each receiver.
The sender's signature could encode parts of its observation most relevant to other
agents' goals (\eg it would be futile to convey coordinates in the signature),
and the message value could contain the agent's own location. For example, in ~\Figref{fig:shapes_3x3},
at $t=6$, we see that when agent 2 passes by blue, agent 4 starts attending to agent 2. Here,
agent 2's signature likely encodes the color it observes (which is blue), and agent 4's query
encodes its goal (which is also blue) leading to high attention probability.
Agent 2's message value encodes coordinates agent 4 has to navigate to,
which it ends up reaching by $t=21$.

SHAPES serves as a flexible testbed for carefully controlling and analyzing
the effect of changing the size of the environment, no. of agents, goal configurations, \etc.
\Figref{fig:shapes_3x3} visualizes learned protocols, and \tableref{table:shapesnumbers}
reports quantitative evaluation for three different configurations --
1) 4 agents, all tasked with finding \texttt{red} in $30 \times 30$ images,
2) 4 agents, all tasked with finding \texttt{red} in $50 \times 50$ images,
3) 4 agents, tasked with finding $[$\texttt{red,red,green,blue}$]$ respectively in $50 \times 50$ images.
We compare TarMAC against two baselines -- 1) without communication, and 2) with communication
but where broadcasted messages are averaged instead of attention-weighted, so
all agents receive the same message vector, similar to~\citet{sukhbaatar_nips16}.
Benefits of communication and attention increase with task complexity
({\scriptsize $30 \times 30 \rightarrow 50 \times 50$ \& find$[$\texttt{red}$]$ $\rightarrow$ find$[$\texttt{red,red,green,blue}$]$}).

\vspace{-5pt}
\subsection{Traffic Junction}
\label{sec:tj}

\textbf{Environment and Task.} The simulated traffic junction environments from~\cite{sukhbaatar_nips16} consist of
cars moving along pre-assigned, potentially intersecting routes on one or more road junctions. The total number of cars is fixed at $N_\textnormal{max}$,
and at every timestep new cars get added to the environment with probability $p_\textnormal{arrive}$. Once a car completes its route, it becomes available to be sampled and added
back to the environment with a different route assignment.
Each car has a limited visibility of a $3 \times 3$ region around it, but is free
to communicate with all other cars. The action space for each car
at every timestep is \texttt{gas} and \texttt{brake}, and the reward consists of a linear time penalty
$-0.01\tau$, where $\tau$ is the number of timesteps since car has been active,
and a collision penalty $r_\textnormal{collision} = -10$.

\begin{table}[h]
\centering
    \resizebox{\columnwidth}{!}{
    \begin{tabular}{lrr}
    \toprule
       & \multicolumn{1}{c}{Easy} & \multicolumn{1}{c}{Hard}    \\ \midrule
        No communication
            & $84.9${\scriptsize $\pm 4.3\%$}  & $74.1${\scriptsize $\pm 3.9\%$}  \\[0.05in]
        CommNet~\citep{sukhbaatar_nips16}
            & $99.7${\scriptsize $\pm 0.1\%$}  & $78.9${\scriptsize $\pm 3.4\%$}  \\[0.05in]
        TarMAC $1$-round
            & $\mathbf{99.9}${\scriptsize $\pm 0.1\%$}   & $84.6${\scriptsize $\pm 3.2\%$}  \\[0.05in]
        TarMAC $2$-round
            & $\mathbf{99.9}${\scriptsize $\pm 0.1\%$} & $\mathbf{97.1}${\scriptsize $\pm 1.6\%$}  \\ \bottomrule
    \end{tabular}} \\[0.05in]
    \caption{Success rates on traffic junction. Our targeted
    $2$-round communication architecture gets a success rate of $97.1${\scriptsize $\pm 1.6\%$}
    on the `hard' variant, significantly outperforming~\citet{sukhbaatar_nips16}.
    Note that $1$- and $2$-round refer to the number of rounds of communication between actions (\Eqref{eq:nextstate}).}
    \label{table:traffic}
    \vspace{-5pt}
\end{table}

\textbf{Quantitative Results.} We compare our approach with CommNet~\citep{sukhbaatar_nips16} on the \texttt{easy}
and \texttt{hard} difficulties of the traffic junction environment. The \texttt{easy} task has one junction of two one-way roads on a $7 \times 7$ grid with $N_\textnormal{max}=5$ and
$p_\textnormal{arrive}=0.30$, while the \texttt{hard} task has four connected
junctions of two-way roads on a $18 \times 18$ grid with $N_\textnormal{max}=20$ and
$p_\textnormal{arrive}=0.05$.
See \Figref{fig:brake}, \ref{fig:attention} for an example
of the four two-way junctions in the \texttt{hard} task. As shown in \tableref{table:traffic},
a no communication baseline has success rates of $84.9${\scriptsize $\pm 4.3\%$}
and $74.1${\scriptsize $\pm 3.9\%$} on \texttt{easy} and \texttt{hard}
respectively. On \texttt{easy}, both CommNet and TarMAC get close to $100\%$.
On \texttt{hard}, TarMAC with $1$-round significantly outperforms
CommNet with a success rate of $84.6${\scriptsize $\pm 3.2\%$}, while $2$-round
further improves on this at $97.1${\scriptsize $\pm 1.6\%$},
which is an $\sim$$18\%$ absolute improvement over CommNet.
We did not see gains going beyond $2$ rounds in this environment.

\textbf{Message size \vs multi-round communication.}
We study performance of TarMAC with varying message value size and number of rounds
of communication on the \texttt{hard} variant of the traffic junction task. As can be seen in~\Figref{fig:msg_sz_rounds},
multiple rounds of communication leads to significantly higher performance than
simply increasing message size, demonstrating the advantage of multi-round communication.
In fact, decreasing message size to a single scalar performs almost as well as $64$-d,
perhaps because even a single real number can be sufficiently partitioned to cover the space of meanings/messages that need to be conveyed.

\begin{figure}[h]
    \centering
    \includegraphics[width=\columnwidth]{figures/main/message_size_vs_round.png}
    \caption{Success rates for $1$ \vs $2$-round \vs message
    size on \texttt{hard}. Performance does not
    decrease significantly even when the message vector is a single scalar, and
    $2$-round communication before taking an
    action leads to significant improvements over $1$-round.}
    \label{fig:msg_sz_rounds}
\end{figure}

\begin{figure*}[t]
    \begin{minipage}[b]{0.28\textwidth}
        \begin{subfigure}{\textwidth}
            \centering
            \includegraphics[width=0.9\textwidth]{figures/main/brake_probs_v2.png}
            \vspace{5pt}
            \caption{Brake probabilities at different locations on the \texttt{hard} traffic
                junction environment. Cars tend to brake close to or right before entering junctions.}
            \label{fig:brake}
        \end{subfigure}
    \end{minipage}\hfill
    \begin{minipage}[b]{0.27\textwidth}
        \begin{subfigure}{\textwidth}
            \centering
            \includegraphics[width=0.9\textwidth]{figures/main/att_probs_v2.png}
            \vspace{5pt}
            \caption{Attention probabilities at different locations.
            Cars are most attended to in the `internal grid' -- right after the $1$st
            junction and before the $2$nd.}
            \label{fig:attention}
        \end{subfigure}
    \end{minipage}\hfill
    \begin{subfigure}{0.41\textwidth}
        \centering
        \vspace{15pt}
        \includegraphics[width=0.9\textwidth]{figures/main/cars_in_sys_vs_att_v2.png}
        \vspace{10pt}
        \caption{No.~of cars being attended to
        1) is positively correlated with total cars,
        indicating that TarMAC is adaptive to dynamic team sizes,
        and 2) is slightly right-shifted, since it takes few steps
        of communication to adapt.}
        \label{fig:carcount}
    \end{subfigure}
    \vspace{5pt}
    \caption{Interpretation of model predictions from TarMAC in the traffic junction environment.}
    \label{fig:traffic}
    \vspace{-10pt}
\end{figure*}

\begin{figure*}[h!]
    \centering
    \includegraphics[width=\textwidth]{figures/main/house3d_nav.jpg}
    \caption{Agents navigating to the \texttt{fireplace} in House3D (marked in yellow).
    Note in particular that agent $4$ is spawned facing away from it.
    It communicates with others, turns to face the \texttt{fireplace}, and moves towards it.}
    \label{fig:house3dnav}
    \vspace{-15pt}
\end{figure*}

\vspace{5pt}
\textbf{Model Interpretation.} Interpreting the learned policies of TarMAC,
\Figref{fig:brake} shows braking probabilities at different locations --
cars tend to brake close to or right before entering traffic junctions, which is reasonable
since junctions have the highest chances for collisions.
Turning our attention to attention probabilities (\Figref{fig:attention}),
we can see that cars are most-attended to when in the `internal grid' -- right after
crossing the $1$st junction and before hitting the $2$nd junction. These attention
probabilities are intuitive -- cars learn to attend to specific sensitive locations
with the most relevant local observations to avoid collisions.
Finally,~\Figref{fig:carcount}
compares total number of cars in the environment
\vs number of cars
being attended to with probability $> 0.1$ at any time. Interestingly, these are (loosely)
positively correlated, with Spearman's $\sigma=0.49$, which shows that TarMAC
is able to adapt to variable number of agents. Crucially, agents learn this dynamic targeting behavior purely from task rewards with
no hand-coding! Note that the right shift between the two curves is
expected, as it takes a few timesteps of communication for team size changes to propagate.
At a relative time shift of $3$, the Spearman's rank correlation between the two curves goes up to $0.53$.

\vspace{-5pt}
\subsection{House3D}
\label{sec:h3d}

Next, we benchmark TarMAC on a cooperative point-goal
navigation task in House3D~\citep{house3d}. House3D provides a rich and diverse
set of publicly-available\footnote{\href{https://github.com/facebookresearch/house3d}{\texttt{github.com/facebookresearch/house3d}}}
$3$D indoor environments, wherein agents do not have access to the top-down map
and must navigate purely from first-person vision. Similar to SHAPES, the agents are tasked
with finding a specified goal (such as `fireplace') within a max no.~of steps,
spawned at random locations in the environment and allowed to communicate and move around.
Each agent gets a shaped reward based on progress towards the specified target.
An episode is successful if all agents end within $0.5$m of the target object
in $500$ navigation steps.

\tableref{table:house3d} shows success rates on a \texttt{find[fireplace]} task
in House3D. A no-communication navigation policy trained with the same
reward structure gets a success rate of $62.1${\scriptsize $\pm 5.3\%$}.
Mean-pooled communication (no attention) performs slightly better with a
success rate of $64.3${\scriptsize $\pm 2.3\%$}, and TarMAC
achieves the best success rate at $68.9${\scriptsize $\pm 1.1\%$}.
TarMAC agents take $82.5$ steps to reach the target on average vs. $101.3$
for no attention vs. $186.5$ for no communication.
\Figref{fig:house3dnav} visualizes a predicted navigation trajectory of $4$ agents.
Note that the communication vectors are significantly more compact ($32$-d)
than the high-dimensional observation space ($224 \times 224$ image),
making our approach particularly attractive for scaling to large agent teams.

\vspace{10pt}
\begin{table}[h]
\centering
\begin{tabular}{lrr}
\toprule
    & \multicolumn{1}{c}{Success rate} & \multicolumn{1}{c}{Avg. \# steps}    \\ \midrule
    No communication
        & $62.1${\scriptsize $\pm 5.3\%$}
        & $186.5$  \\[0.05in]
    No attention
        & $64.3${\scriptsize $\pm 2.3\%$}
        & $101.3$ \\[0.05in]
    TarMAC
        & $\mathbf{68.9}${\scriptsize $\pm 1.1\%$}
        & $\mathbf{82.5}$  \\ \bottomrule
\end{tabular} \\[0.09in]
\caption{$4$-agent \texttt{find[fireplace]} navigation task in House3D.}
\label{table:house3d}
\end{table}

Note that House3D is a challenging testbed for multi-agent reinforcement
learning. To get to $\sim$$100\%$ accuracy, agents have to deal with high-dimensional
visual observations, be able to navigate long action sequences (up to $\sim$$500$ steps),
and avoid getting stuck against objects, doors, and walls.

\vspace{-5pt}
\subsection{Mixed and Competitive Environments}
\label{sec:competitive}

\begin{figure*}[t]
    \begin{minipage}[b]{0.45\textwidth}
        \centering
        \begin{subfigure}{1.0\textwidth}
            \centering
            \includegraphics[width=0.9\textwidth]{figures/main/pp_easy_steps.png}
            \vspace{5pt}
            \caption{$3$ agents, $5\times5$ grid,
                vision=$0$, max steps=$20$}
            \label{fig:pp_easy}
            \vspace{10pt}
        \end{subfigure}
    \end{minipage}\hfill
    \begin{minipage}[b]{0.45\textwidth}
        \centering
        \begin{subfigure}{\textwidth}
            \centering
            \includegraphics[width=0.9\textwidth]{figures/main/pp_hard_steps.png}
            \vspace{5pt}
            \caption{$10$ agents, $20\times20$ grid,
                vision=$1$, max steps=$80$}
            \label{fig:pp_hard}
            \vspace{10pt}
        \end{subfigure}
    \end{minipage}\hfill
    \vspace{-5pt}
    \caption{Average no.~of steps to complete an episode (lower is better)
        during training in the Predator-Prey mixed environment.
        IC3Net + TarMAC converges much faster than IC3Net,
        demonstrating that attentional communication helps.
        Shaded region shows $95\%$ CI.}
    \label{fig:pp}
    \vspace{-10pt}
\end{figure*}

\begin{table*}[t]
\centering
\begin{tabular}{lccc}
    & \multicolumn{1}{c}{{\scriptsize $3$ agents, $5 \times 5$,}}
        & \multicolumn{1}{c}{{\scriptsize $5$ agents, $10 \times 10$,}}
        & \multicolumn{1}{c}{{\scriptsize $10$ agents, $20 \times 20$,}} \\[-0.05in]
    & \multicolumn{1}{c}{{\scriptsize vision=$0$, max steps=$20$}}
        & \multicolumn{1}{c}{{\scriptsize vision=$1$, max steps=$40$}}
        & \multicolumn{1}{c}{{\scriptsize vision=$1$, max steps=$80$}}    \\ \midrule
    {CommNet~\cite{sukhbaatar_nips16}}
        & {$9.1$}{\scriptsize $\pm 0.1$} & {$13.1$}{\scriptsize $\pm 0.01$} & {$76.5$}{\scriptsize $\pm 1.3$}  \\ [0.05in]
    {IC3Net~\cite{singh_iclr19}}
        & {$8.9$}{\scriptsize $\pm 0.02$} & {$13.0$}{\scriptsize $\pm 0.02$} & {$52.4$}{\scriptsize $\pm 3.4$} \\ [0.05in]
    {IC3Net $+$ TarMAC}
        & {$8.31$}{\scriptsize $\pm 0.06$} & {$\mathbf{12.74}$}{\scriptsize $\pm 0.08$} & {$41.67$}{\scriptsize $\pm 5.82$} \\ [0.05in]
    {IC3Net $+$ TarMAC (2-round)}
        & {$\mathbf{7.24}$}{\scriptsize $\pm 0.08$} & -- & {$\mathbf{35.57}$}{\scriptsize $\pm 3.96$} \\ \bottomrule
\end{tabular} \\[0.1in]
\caption{Average number of steps taken to complete an episode (lower is better)
        at convergence in the Predator-Prey mixed environment.}
\label{table:ppnumbers}
\vspace{-15pt}
\end{table*}

Finally, we look at how to extend TarMAC to
mixed and competitive scenarios. Communication via sender-receiver soft
attention in TarMAC is poorly suited for competitive scenarios,
since there is always ``leakage'' of the agent's state as a message to other
agents via a low but non-zero attention probability, thus compromising its
strategy and chances of success. Instead, an agent should first be able to independently decide if it wants to
communicate at all, and then direct its message to specific recipients if it does.

The recently proposed IC3Net architecture by~\citet{singh_iclr19} addresses
the former -- learning when to communicate. At every timestep, each agent in IC3Net
predicts a hard gating action to decide if it wants to communicate.
At the receiving end, messages from agents who decide to communicate are
averaged to be the next input message.
Replacing this message averaging with our sender-receiver soft attention,
while keeping the rest of the architecture and training details the same as
IC3Net,
should provide an inductive bias for more flexible communication strategies,
since this model (IC3Net + TarMAC) can learn both when to
communicate and whom to address messages to.

We evaluate IC3Net + TarMAC on the Predator-Prey environment from~\citet{singh_iclr19},
consisting of $n$ predators, with limited vision, moving around
(with a penalty of $r_{\text{explore}}=-0.05$ per timestep)
in search of a stationary prey. Once a predator reaches a prey, it keeps getting positive
reward $r_{\text{prey}}=0.05$ till end of episode \ie till other agents reach prey or
maximum no.~of steps. The prey gets $0.05$ per timestep only till the first
predator reaches it, so it has incentive to not communicate its location.
We compare average no. of steps for
agents to reach the prey during training (\Figref{fig:pp}) and at convergence (\tableref{table:ppnumbers}).
\Figref{fig:pp} shows that using TarMAC with IC3Net leads to significantly faster convergence
than IC3Net alone, and \tableref{table:ppnumbers} shows that TarMAC agents
reach the prey faster.
Results are averaged over $3$ independent runs with different seeds.



